\centerline{\bf \Huge Домашнее задание.}
\centerline{\bf \Huge Метод Штурма}

\begin{flushright}
        \scriptsize{...шесть или пять полов напридумывали! \\Трансформеры, транс... Я даже не понимаю, что это такое!\\ВВП}
\end{flushright}

\problem{1!}\textbf{(Повторение)}
Пусть сумма двух положительных чисел a и b фиксирована. Как изменяются следующие выражения при <<сближении>> a и b:

\begin{enumerate}
    \item  $ab$
    \item  $a^2 + b^2;$
    \item  $\dfrac{1}{a} + \dfrac{1}{b};$
    \item  $\sqrt{a} + \sqrt{b};$
    \item  $a^n + b^n;$
    \item  $\dfrac{1}{a^n} + \dfrac{1}{b^n}.$
\end{enumerate}

\problem{2!}
У Яны есть палочка длины 1 метр. Она хочет распилить ее на 12 палочек и сколотить из них каркас прямоугольного параллелепипеда. Какой наибольший объем может иметь данный параллелепипед?

\problem{3!}
Вершины графа со $100$ вершинами раскрашены в $7$ цветов так, что концы каждого ребра --- разного цвета. Какое наибольшее число ребер может быть в этом графе?

\problem{4}
\textbf{(Неравенство Гюйгенса)} Пусть $a_1, a_2, \ldots,a_n$ положительные числа, $t$ их среднее геометрическое. Докажите неравенство $(1+a_1)(1+a_2)\ldots(1+a_n) \geq t^n$

\problem{5}
\text { Положительные числа } $a_1, a_2, \ldots, a_n$ \text{, меньше 1. Докажите, что } \[\left[a_1+\ldots+a_n\right]+\left\{a_1+\ldots+a_n\right\}^2 \geq a_1^2+\ldots+a_n^2\]

\problem{6}
а) Известно, что $x_1>x_2$ и $y_1>y_2.$ Что больше: $x_1y_1+x_2y_2$ или $x_1y_2+x_2y_1?$

б) Докажите транс-неравенство: если $a_1 \geq a_2 \geq \ldots a_n \geq 0, b_1 \geq b_2 \geq \ldots \geq b_n \geq 0$ и $c_1, c_2, \ldots, c_n$ --- некоторая перестановка чисел $b_1, b_2, \ldots, b_n,$ то
\[a_1b_1+a_2b_2+ \ldots +a_nb_n \leq a_1c_1 + a_2c_2 + \ldots + a_nc_n \leq a_1b_n + a_2b_{n-1} + \ldots + a_nb_1.\]